\section{Variance bounds and dependent-sampling proofs}
\label{sec:annex-p15-variance-dependent-concentration}

\begin{annexitem}{annex:p15.upper-moment-extremal}{thm:p15-upper-moment-extremal}{A quadratic majorant on the entire upper half-line}
\begin{proof}
The case \(v=0\) follows from \(EX^2=0\). Suppose \(v>0\), put
\(a=-v/b<b\), and let \(F(x)=e^{hx}\). There is a unique quadratic
\[
 Q(x)=F(a)+F'(a)(x-a)
 +\frac{F(b)-F(a)-F'(a)(b-a)}{(b-a)^2}(x-a)^2
\]
matching \(F,F'\) at \(a\) and \(F\) at \(b\).
For \(x\ne a,b\), subtract
\(K(z-a)^2(z-b)\) from \(F(z)-Q(z)\), choosing \(K\) to make the
difference zero also at \(z=x\). The resulting function has a double
zero at \(a\), and zeros at \(b,x\).
Repeated Rolle's theorem on the interval containing these points gives
a zero of its third derivative: between distinct zeros the derivative
has a zero, and the double zero contributes the additional zero of the
first derivative at \(a\). Repeating this count twice more gives
\[
 F(x)-Q(x)=\frac{F'''(\xi)}6(x-a)^2(x-b)
\]
for some intervening \(\xi\). Since \(F'''(\xi)=h^3e^{h\xi}>0\),
this difference is nonpositive for every \(x\leq b\).
At \(x=a,b\) the same assertion holds by equality.
This proves the majorant on all of \((-\infty,b]\), not just between
the two interpolation nodes.

Let \(Z\) have the two-point distribution stated in the body.
Direct calculation gives \(EZ=0\) and \(EZ^2=v\):
the first-moment numerator is \(-bv+bv=0\), and the second is
\((v^2+vb^2)/(b^2+v)=v\).
Expectations of the quadratic \(Q\) agree for \(X\) and \(Z\), because
their first two moments agree. They are finite.
Therefore
\[
 Ee^{hX}\leq EQ(X)=EQ(Z)=Ee^{hZ}=M_{b,v}(h).
\]
The two-point law attains equality, completing the editorial proof
of the moment lemma without importing an unread proof from Bennett.
\end{proof}
\end{annexitem}

\begin{annexitem}{annex:p15.variance-moment-concavity}{thm:p15-variance-moment-concavity}{The source's concavity calculation, including zero variance}
\begin{proof}
Write \(y=1+u\) and
\[
 F(y)=1-y^{-1}+y^{-1}e^{-cy},\qquad f_c(u)=c+\log F(y).
\]
For \(y\geq1\), \(F(y)>0\), including \(F(1)=e^{-c}>0\).
Differentiation gives
\[
 F'(y)=y^{-2}\{1-(1+cy)e^{-cy}\}>0,
\]
and
\[
 F''(y)=-2y^{-3}e^{-cy}
       \{e^{cy}-1-cy-(cy)^2/2\}<0.
\]
The strict signs follow from the positive remainder terms of the
exponential series for \(cy>0\), or its integral Taylor remainder.
Hence
\[
 f_c'(u)=F'(y)/F(y)>0,\qquad
 f_c''(u)=\frac{F(y)F''(y)-F'(y)^2}{F(y)^2}<0.
\]
The same formulas hold at \(u=0\) as right derivatives. Substitution
gives \(f_c(0)=0\) and \(f_c'(0)=e^c-1-c\).

Independence and the moment lemma yield
\[
 \log E e^{h\sum_iX_i}
 \leq\sum_i f_{bh}(v_i/b^2)
 \leq n f_{bh}(\sigma^2/b^2).
\]
The second inequality is finite Jensen for the concave function just
proved; it permits some \(v_i=0\).
Exponentiation gives the first asserted moment bound.
Concavity also gives
\(f_c(u)\leq f_c(0)+u f_c'(0)=u(e^c-1-c)\).
Substitute \(c=bh,u=\sigma^2/b^2\) to obtain the second bound.
\end{proof}
\end{annexitem}

\begin{annexitem}{annex:p15.variance-tail-chain}{thm:p15-variance-tail-chain}{Two optimizations and the complete simpler-bound comparisons}
\begin{proof}
The indicator of \(\sum_iX_i\geq nt\) is at most
\(e^{h(\sum_iX_i-nt)}\) for \(h>0\).
By the preceding lemma its expectation is bounded by
\[
 \left\{
 \frac{b^2}{b^2+\sigma^2}e^{-h(t+\sigma^2/b)}
 +\frac{\sigma^2}{b^2+\sigma^2}e^{h(b-t)}
 \right\}^{n}.
\]
For \(0<t<b\), differentiate the positive expression in braces.
Its derivative vanishes exactly when
\[
 h=h_*=\frac{b}{b^2+\sigma^2}
       \log\frac{1+bt/\sigma^2}{1-t/b}>0 .
\]
The second derivative is strictly positive, so this point is its
global minimum on \(h>0\).
For an explicit substitution, set
\(p=\sigma^2/(b^2+\sigma^2)\) and
\(q=(\sigma^2+bt)/(b^2+\sigma^2)\). The expression is the centered
two-point Bernoulli exponential expression, whose value at \(h_*\) is
\((p/q)^q\{(1-p)/(1-q)\}^{1-q}\).
Here \(q/p=1+bt/\sigma^2\) and
\((1-q)/(1-p)=1-t/b\); substitution gives \(B_1\) exactly.

The tangent moment bound instead gives the following route to
\(B_2\), the alternative direct proof of Bennett's bound that Hoeffding
indicates \cite[p.~24]{Hoeffding1963BoundedSums}:
\[
 \Pr\{\overline X\geq t\}
 \leq \exp\{-nht+n\sigma^2(e^{hb}-1-hb)/b^2\}.
\]
Its exponent has derivative
\(-nt+n\sigma^2(e^{hb}-1)/b\), and strictly positive second derivative.
The minimum is at \(hb=\log(1+bt/\sigma^2)\).
The minimized exponent is
\[
 -\frac{n\sigma^2}{b^2}
       \{(1+\lambda)\log(1+\lambda)-\lambda\}
 =-\tau h_1(\lambda).
\]
This optimization is valid for every \(t>0\).
Pointwise, the first moment envelope is at most the tangent envelope;
their infima over \(h>0\) therefore obey \(B_1\leq B_2\) on the
interior domain.

To prove the Bernstein comparison put
\(H(\lambda)=(1+\lambda)\log(1+\lambda)-\lambda\) and
\(R(\lambda)=3\lambda^2/\{2(\lambda+3)\}\).
Both their values and first derivatives vanish at zero, and
\[
 H''(\lambda)-R''(\lambda)
 =\frac1{1+\lambda}-\frac{27}{(\lambda+3)^3}\geq0.
\]
Indeed \((\lambda+3)^3-27(1+\lambda)=\lambda^3+9\lambda^2\geq0\).
Integrating twice gives \(H\geq R\), or \(h_1\geq h_2\).
Since \(\tau>0\), this is \(e^{-\tau h_1(\lambda)}\leq e^{-\tau h_2(\lambda)}\):
the bound (2.12) is never larger than Bernstein's (2.13), which is the
comparison Hoeffding attributes to Bennett
\cite[p.~17]{Hoeffding1963BoundedSums}.
For the Prohorov comparison, differentiation shows
\[
 \log(1+\lambda)\geq\frac{2\lambda}{2+\lambda},
\]
since the derivative of the difference is
\(\lambda^2/\{(1+\lambda)(2+\lambda)^2\}\geq0\).
It follows that
\[
 h_1-\tfrac12\log(1+\lambda)
 =(\tfrac12+\lambda^{-1})\log(1+\lambda)-1\geq0.
\]
Also \(\sqrt{\lambda^2+4}\leq\lambda+2\), whence
\(\operatorname{arsinh}(\lambda/2)
 =\log\{(\lambda+\sqrt{\lambda^2+4})/2\}\leq\log(1+\lambda)\).
Thus \(h_1\geq h_3\). Substituting \(\lambda,\tau\) into the
Bernstein expression gives exactly the displayed denominator.

For \(t\uparrow b\), the events \(\{\overline X\geq t\}\) decrease
to \(\{\overline X\geq b\}\). Probability is continuous on such a
decreasing sequence. The factor involving
\((1-t/b)\log(1-t/b)\) tends to zero in the log expression, leaving
\(\{\sigma^2/(b^2+\sigma^2)\}^n\). Since \(\overline X\leq b\),
larger thresholds have probability zero.
If \(\sigma^2=0\), all the nonnegative \(v_i\) vanish and every \(X_i\)
is zero almost surely. If \(b=0\), \(-X_i\) is nonnegative with mean zero,
so again it vanishes. For \(b<0\), \(EX_i\leq b<0\) contradicts
centering. Finally positive upper and lower deviations are disjoint;
apply the proved upper-tail statements separately to \(X_i\) and
\(-X_i\) when both upper bounds are available.
\end{proof}
\end{annexitem}

\begin{annexitem}{annex:p15.variance-maximal}{thm:p15-variance-maximal}{A first-crossing proof with the terminal exponential moment}
\begin{proof}
Let \(A_j\) be the event that \(j\) is the first index with
\(S_j\geq nt\), for \(1\leq j\leq n\). These are disjoint events
depending only on \(X_1,\ldots,X_j\), and their union is the maximum
event, since \(S_0=0<nt\).
Every positive exponential moment is finite because \(X_i\leq b\).
For each future increment, \(e^{hx}\geq1+hx\) and \(EX_i=0\)
give \(Ee^{hX_i}\geq1\). Independence of the future from the past yields
\[
 E[{\bf1}_{A_j}e^{hS_n}]
 =E[{\bf1}_{A_j}e^{hS_j}]\prod_{i=j+1}^nEe^{hX_i}
 \geq e^{hnt}\Pr(A_j).
\]
For \(j=n\) the empty product is one.
Sum over the disjoint events and bound the sum of expectations by
\(Ee^{hS_n}\). This proves the stated inequality, with exactly the
terminal moment already bounded and optimized in the variance theorem.
If \(t>b\), every \(S_j\leq jb\leq nb<nt\).
At \(t=b>0\), a crossing is possible only at \(j=n\), so the earlier
endpoint bound applies directly. If the total variance is zero, every
increment vanishes almost surely.

For Kolmogorov's inequality, let now \(A_j\) be the event that \(j\) is
the first index with \(|S_j|\geq x\), \(1\leq j\leq n\). These events are
disjoint, each depends only on \(X_1,\ldots,X_j\), and their union is
\(\{\max_{1\leq j\leq n}|S_j|\geq x\}\). Put \(R_j=S_n-S_j\), a function
of \(X_{j+1},\ldots,X_n\) with \(ER_j=0\) and \(ER_j^2<\infty\). The
variable \(S_j{\bf1}_{A_j}\) is square integrable and independent of
\(R_j\), so \(E[S_j{\bf1}_{A_j}R_j]=E[S_j{\bf1}_{A_j}]\,ER_j=0\). Hence
\[
 E[S_n^2{\bf1}_{A_j}]
 =E[S_j^2{\bf1}_{A_j}]+2E[S_j{\bf1}_{A_j}R_j]+E[R_j^2{\bf1}_{A_j}]
 \geq E[S_j^2{\bf1}_{A_j}]\geq x^2\Pr(A_j),
\]
the last step because \(S_j^2\geq x^2\) on \(A_j\). Summing over the
disjoint \(A_j\) gives
\(ES_n^2\geq\sum_jE[S_n^2{\bf1}_{A_j}]\geq x^2\Pr\{\max_j|S_j|\geq x\}\).
The structure is that of the exponential argument above: a first-crossing
decomposition, followed by the observation that the future increments do
not lower the chosen moment on average. No upper bound on the \(X_i\) was
used here.
\end{proof}
\end{annexitem}

\begin{annexitem}{annex:p15.variance-bound-comparison}{thm:p15-variance-bound-comparison}{Exact log-ratio algebra and variance monotonicity}
\begin{proof}
Put \(r=\sigma^2/b^2\), \(w=t/b\). The negative log bounds divided by
\(n\) are
\[
 D=\frac{(r+w)\log(1+w/r)+(1-w)\log(1-w)}{1+r},
 \qquad E=(r+w)\log(1+w/r)-w.
\]
Thus \(B_1/B_2=e^{-n(D-E)}\).
With \(v=w/(r+w)\), let \(K=v+w-vw>0\), so
\(r=w(1-v)/v\) and \(1+r=K/v\). Directly collecting the logarithms gives
\[
 D-E=\frac{
  (w^2/v)(1-v)\log(1-v)+v(1-w)\log(1-w)+wK}{K}
 =\frac{vw^2}{K}\{\rho(v)+\rho(w)\}.
\]
Since \(K/(vw)=v^{-1}+w^{-1}-1\), this is the asserted formula.

Integrating the convergent power series for \(-\log(1-x)\), or
multiplying it by \(1-x\), gives
\((1-x)\log(1-x)+x=\sum_{\ell\geq2}x^\ell/\{\ell(\ell-1)\}\).
Subtract its quadratic term and divide by \(x^2\) to obtain the
series for \(\rho\). Every coefficient is positive.
For \(0<x<1\), the first term gives \(\rho(x)>x/6\).
Moreover
\(\sum_{j\geq1}1/\{(j+1)(j+2)\}=1/2\), by telescoping, so
\(\rho(x)<x/2\), strictly because \(x^j<x\) for \(j>1\).
Substitution gives the first two bounds for \(\phi\).
For the last, the middle upper bound equals
\(vw^2(v+w)/(2K)\), whereas
\(t^3/(2b\sigma^2)=w^3/(2r)=vw^2/\{2(1-v)\}\).
The required strict comparison follows from
\((v+w)(1-v)<v+w-vw=K\).

For fixed \(b,\sigma^2\), both \(v,w\) tend to zero linearly with \(t\);
\(\rho(v)+\rho(w)=(v+w)/6+O(t^2)\) and
\(K=v+w+O(t^2)\). Consequently
\(\phi=vw^2/6+O(t^4)=t^3/(6b\sigma^2)+O(t^4)\).

For the small-\(\lambda\) regime, note first that
\(t^3/(6b\sigma^2)=\lambda w^2/6\) and \(v=\lambda/(1+\lambda)<\lambda\).
By the formula for \(\phi\) and \(K/(vw)=v^{-1}+w^{-1}-1\),
\[
 \frac{\phi(v,w)}{\lambda w^2/6}
 =\frac v\lambda\cdot\frac{6\{\rho(v)+\rho(w)\}}{K},\qquad
 K=v+w-vw .
\]
Lower bound: \(\rho(x)>x/6\) and \(K<v+w\) make the second factor
larger than one, and the first factor is \(v/\lambda=1/(1+\lambda)\).
Upper bound: let \(m=\max\{v,w\}\leq M<1\). In the series for
\(\rho\), every coefficient after the first is at most
\(1/\{(3)(4)\}=1/12\), so for \(0<x\leq m\)
\[
 \rho(x)\leq\frac x6+\frac1{12}\sum_{j\geq2}x^j
 =\frac x6+\frac{x^2}{12(1-x)}
 \leq x\Bigl(\frac16+\frac m{12(1-m)}\Bigr).
\]
It follows that
\(6\{\rho(v)+\rho(w)\}\leq(v+w)(2-m)/\{2(1-m)\}\). Also
\(vw\leq m\min\{v,w\}\leq m(v+w)\), so \(K\geq(v+w)(1-m)\). With
\(v/\lambda<1\) this gives
\[
 \frac{\phi(v,w)}{\lambda w^2/6}<\frac{2-m}{2(1-m)^2}
 \leq\frac{2-M}{2(1-M)^2},
\]
the last step because \(m\mapsto\log(2-m)-2\log(1-m)\) has derivative
\(-1/(2-m)+2/(1-m)>0\) on \([0,1)\).
Next, \(w=t/b=(bt/\sigma^2)(\sigma^2/b^2)=\lambda\sigma^2/b^2\), so
\(\sigma^2\leq b^2\) gives \(w\leq\lambda\) and \(M=\lambda\). Under
\(a\leq X_i\leq b\) the bound \(\sigma^2\leq-ab\) is proved below, and
\(-ab\leq b^2\) when \(-a\leq b\).
For the fixed-\(w\) limit, \(\lambda\downarrow0\) forces \(v\downarrow0\),
\(v/\lambda=1/(1+\lambda)\to1\), \(0<\rho(v)<v/2\to0\) and \(K\to w\).
Hence the displayed quotient tends to \(6\rho(w)/w\). From the series,
\[
 \frac{6\rho(w)}w=6\sum_{j\geq1}\frac{w^{j-1}}{(j+1)(j+2)}
 =1+\frac w2+6\sum_{j\geq3}\frac{w^{j-1}}{(j+1)(j+2)}>1+\frac w2 .
\]
The product of this limit with \(t^3/(6b\sigma^2)\) is
\[
 \frac{t^3}{6b\sigma^2}\cdot\frac{6b}t\,\rho(w)
 =\frac{t^2}{\sigma^2}\,\rho(t/b).
\]
At \(w=1/2\),
\(\rho(1/2)=4\{\tfrac12\log\tfrac12+\tfrac12-\tfrac18\}=\tfrac32-2\log2\),
and \(6\rho(1/2)/(1/2)=12\rho(1/2)=18-24\log2\approx1.3645\).

For the expansions of \(h_1,h_2\), let \(0<\lambda\leq1\). The series
\(\log(1+\lambda)=\sum_{k\geq1}(-1)^{k+1}\lambda^k/k\) converges, and
\[
 \frac{\log(1+\lambda)}\lambda
 =\sum_{k\geq1}\frac{(-1)^{k+1}\lambda^{k-1}}k
 =1+\sum_{k\geq1}\frac{(-1)^k\lambda^k}{k+1}.
\]
Adding the two series term by term and subtracting one gives
\[
\begin{split}
 h_1(\lambda)&=\log(1+\lambda)+\frac{\log(1+\lambda)}\lambda-1\\
 &=\sum_{k\geq1}(-1)^{k+1}\lambda^k\Bigl(\frac1k-\frac1{k+1}\Bigr)
 =\sum_{k\geq1}\frac{(-1)^{k+1}\lambda^k}{k(k+1)} .
\end{split}
\]
Its first two terms are \(\lambda/2\) and \(-\lambda^2/6\). The terms
alternate in sign and their absolute values
\(\lambda^k/\{k(k+1)\}\) strictly decrease to zero, because
\(\lambda\leq1\) and \(k(k+1)\) increases. The alternating-series
estimate therefore places the remainder after the second term strictly
between zero and its first term \(\lambda^3/12\). For \(h_2\), direct
computation over the common denominator \(2(1+\lambda/3)\) gives
\[
 \frac\lambda{2(1+\lambda/3)}-\frac\lambda2+\frac{\lambda^2}6
 =\frac{\lambda-\lambda(1+\lambda/3)+(\lambda^2/3)(1+\lambda/3)}
       {2(1+\lambda/3)}
 =\frac{\lambda^3}{18(1+\lambda/3)},
\]
valid for every \(\lambda>0\). Both remainders are \(O(\lambda^3)\) as
\(\lambda\downarrow0\).

For the strict increase of \(B_1\) in \(\sigma^2\), fix \(n\), \(b\) and
\(0<t<b\). For \(v\geq0\) and \(h>0\), put \(c=bh\) and \(u=v/b^2\);
multiplying numerator and denominator by \(b^2\),
\[
 e^{f_c(u)}=\frac{e^{-hv/b}+(v/b^2)e^{hb}}{1+v/b^2}
 =\frac{b^2e^{-hv/b}+ve^{hb}}{b^2+v}=M_{b,v}(h).
\]
By \cref{thm:p15-variance-moment-concavity}, \(f_c'>0\) on \([0,\infty)\),
so \(f_c\) is strictly increasing there, and hence
\(M_{b,v_1}(h)<M_{b,v_2}(h)\) whenever \(0\leq v_1<v_2\) and \(h>0\).
Write \(\Psi_{\sigma^2}(h)=e^{-hnt}M_{b,\sigma^2}(h)^n\), the bound minimized
in the proof of \cref{thm:p15-variance-tail-chain}: there the expression
in braces is \(e^{-ht}M_{b,\sigma^2}(h)\). That proof shows that, for each
\(\sigma^2>0\), \(\Psi_{\sigma^2}\) attains its minimum over \(h>0\) at the
single point \(h_*(\sigma^2)>0\) displayed there, and that the minimum
value is \(B_1(\sigma^2)\). Let \(0<\sigma_1^2<\sigma_2^2\) and
\(h_2=h_*(\sigma_2^2)\). Then
\[
 B_1(\sigma_1^2)\leq\Psi_{\sigma_1^2}(h_2)
 <\Psi_{\sigma_2^2}(h_2)=B_1(\sigma_2^2).
\]
The first step holds because \(B_1(\sigma_1^2)\) is a minimum over all
\(h>0\). The second holds because \(0<M_{b,\sigma_1^2}(h_2)<M_{b,\sigma_2^2}(h_2)\),
the map \(y\mapsto y^n\) is strictly increasing on \((0,\infty)\) for
\(n\geq1\), and \(e^{-h_2nt}>0\). So \(B_1\) is strictly increasing in
\(\sigma^2\), which is the increase Hoeffding asserts for (2.8)
\cite[p.~20]{Hoeffding1963BoundedSums}; in particular it is
nondecreasing, which is all the substitution below uses.
For \(a\leq X_i\leq b\),
\((X_i-a)(X_i-b)\leq0\); taking expectations and using \(EX_i=0\)
gives \(EX_i^2\leq-ab\).
At \(v=-ab\), the two interpolation support values are \(a,b\),
with probabilities \(b/(b-a),-a/(b-a)\).
Its normalized upper-endpoint probability is \(p=-a/(b-a)\);
threshold \(t\) becomes \(q=p+t/(b-a)\).
The formula already computed is exactly
\(e^{-nD(q\Vert p)}\), the first common-range bound.

For \(B_2\), the tail proof minimizes over \(h>0\) the functions
\(\exp\{-nht+n\sigma^2(e^{hb}-1-hb)/b^2\}\). Each is nondecreasing in
\(\sigma^2\) because \(e^y-1-y\geq0\) for every real \(y\); hence their
infimum \(B_2\) is nondecreasing in \(\sigma^2\). With
\(0<\sigma^2\leq-ab\) this gives
\(\Pr\{\overline X\geq t\}\leq B_2(\sigma^2)\leq B_2(-ab)\) for every
\(t>0\); if \(\sigma^2=0\) the probability is zero.
At \(\sigma^2=-ab\), \(\lambda=bt/(-ab)=-t/a\), so
\(1+\lambda^{-1}=1-a/t\) and \(\log(1+\lambda)=\log(1-t/a)\), and
\(\tau h_1(\lambda)\) is the displayed exponent.
For \(0<t<b\), the inequality \(B_1\leq B_2\) holds for every positive
variance parameter, in particular for \(-ab\); combined with the
identification of \(B_1(-ab)\) just proved, the variance-free \(B_2\)
bound is at least the common-range bound.
The exact ratio is the first display of this proof with
\(\sigma^2=-ab\), for which \(v=bt/(bt-ab)=t/(t-a)\).

For the cited calculation, the values \(1\) and \(-1/4\) give
\(-ab=1/4\), which equals the common variance; the bounds are therefore
the \(B_1,B_2\) values computed in
\cref{annex:p15.dependent-concentration-examples}. There
\(p=-a/(b-a)=1/5\) and \(q=p+t/(b-a)=2/5\), and
\(e^{-20D(2/5\Vert1/5)}=2^{-8}(3/4)^{-12}=2^{16}/3^{12}\).
\end{proof}
\end{annexitem}

\begin{annexitem}{annex:p15.variance-optimality-boundary}{thm:p15-variance-optimality-boundary}{Attainable moment envelopes and a different tail optimum}
\begin{proof}
Take independent copies of the moment-extremizing law with \(v=\sigma^2\).
Their joint exponential moment is exactly \(M_{b,\sigma^2}(h)^n\)
for every \(h>0\). The optimization in the tail proof therefore gives
exactly \(B_1\) as their best single-exponential Markov bound.
At \(t=b\), the event requires every variable to equal \(b\); its
probability is the product of the endpoint masses, equal to \(B_1\).

For a general centered \(X\) of variance one and \(t>0\), every
\(a>0\) gives
\[
 \Pr\{X\geq t\}\leq\frac{E(X+a)^2}{(t+a)^2}
                         =\frac{1+a^2}{(t+a)^2}.
\]
This uses only the implication \(X\geq t\Rightarrow(X+a)^2\geq(t+a)^2\).
At \(a=1/t\) the bound is \(1/(1+t^2)\). For \(t=1/2\) it is \(4/5\).
The two-point law at \(1/2,-2\) with probabilities \(4/5,1/5\)
has mean zero, variance one, and upper bound one, and attains that tail.
Substitution into the exponential formula gives \(2/3^{3/4}\).
It exceeds \(4/5\), since the fourth powers satisfy
\(16/27>256/625\). Thus the distinction concerns an actual uniform
tail optimum, not just an example of slack for a chosen law.

For the general exponential majorant, its two hypotheses imply
\({\bf1}_{\{W\geq a\}}\leq Q(W)\) almost surely.
Every exponential is integrable by assumption, so the finite signed
linear combination is integrable. Taking expectations and using
linearity gives the stated bound. If \(c_j<0\), multiplying a moment
upper bound by \(c_j\) reverses its inequality; a lower bound or exact
moment is needed for an upper bound on that summand.

For the final assertion, write \(B_1=e^{-nD}\),
\(\tau=nw\), as in the preceding proof. We have
\(D/w=h_1(\lambda)+\phi(v,w)/w\geq h_1(\lambda)\).
Fix \(\lambda>0\), hence \(v=\lambda/(1+\lambda)\), and let
\(w\downarrow0\), setting \(\sigma^2/b^2=w/\lambda\).
These are admissible parameters. The explicit formula shows
\[
 \frac{\phi(v,w)}w
 =\frac{\rho(v)+\rho(w)}{v^{-1}+w^{-1}-1}\longrightarrow0.
\]
Hence \(\inf_{0<w<1}D/w=h_1(\lambda)\).
Any bound \(B_1\leq e^{-\tau H(\lambda)}\) uniformly in these
parameters requires \(H(\lambda)\leq D/w\) for every \(w\), and
therefore \(H(\lambda)\leq h_1(\lambda)\).
\end{proof}
\end{annexitem}

\begin{annexitem}{annex:p15.variance-upper-bound-crucial}{thm:p15-variance-upper-bound-crucial}{A one-summand witness for the Chebyshev value}
\begin{proof}
\emph{Upper bound.} Let \((X_i)\in\mathcal C(n,\sigma^2)\) and
\(W=\sum_iX_i=n\overline X\). Then \(EW=0\), and independence makes the
cross terms \(E[X_iX_j]=EX_i\,EX_j=0\) for \(i\neq j\), so
\(EW^2=\sum_iEX_i^2=n\sigma^2<\infty\). The one-sided bound of
\cref{thm:p15-conc-benchmarks}, with \(v=n\sigma^2\) and \(u=nt>0\), gives
\[
 \Pr\{\overline X\geq t\}=\Pr\{W\geq nt\}
 \leq\frac{n\sigma^2}{n\sigma^2+n^2t^2}
 =\frac{\sigma^2}{\sigma^2+nt^2}.
\]

\emph{Attainment.} Write \(p=\sigma^2/(\sigma^2+nt^2)\in(0,1)\),
\(s=nt\) and \(c=\sigma^2/t\). The variables \(X_2,\ldots,X_n\) are the
constant zero, so all \(n\) variables are jointly independent, and each
of \(X_2,\ldots,X_n\) has mean and second moment zero. For \(X_1\),
\[
 EX_1=ps-(1-p)c
 =\frac{\sigma^2\,nt-nt^2\,\sigma^2/t}{\sigma^2+nt^2}=0,
\]
\[
 EX_1^2=ps^2+(1-p)c^2
 =\frac{\sigma^2n^2t^2+nt^2\,\sigma^4/t^2}{\sigma^2+nt^2}
 =\frac{n\sigma^2(nt^2+\sigma^2)}{\sigma^2+nt^2}=n\sigma^2 .
\]
So \(n^{-1}\sum_iEX_i^2=\sigma^2\) and the family lies in
\(\mathcal C(n,\sigma^2)\). Since \(\overline X=X_1/n\) and
\(-c<0<s\), the event \(\{\overline X\geq t\}\) is \(\{X_1=nt\}\), of
probability \(p\). The supremum therefore equals \(p\) and is attained.

\emph{No better variance bound.} If \(\beta(n,t,\sigma^2)\) bounds
\(\Pr\{\overline X\geq t\}\) for every member of
\(\mathcal C(n,\sigma^2)\), it bounds the attaining family, so
\(\beta(n,t,\sigma^2)\geq\sigma^2/(\sigma^2+nt^2)\).

\emph{Large upper bounds.} Let \(b\geq nt\). The attaining family takes
values in \(\{nt,-\sigma^2/t\}\) and \(\{0\}\), all at most \(nt\leq b\),
so it satisfies \cref{def:p15-upper-variance-setting} with this \(b\).
That class is contained in \(\mathcal C(n,\sigma^2)\). Its supremum is
therefore at most \(\sigma^2/(\sigma^2+nt^2)\) by the upper bound, and at
least the same number by the witness. For fixed \(b>0\) and \(t>0\),
\(nt>b\) once \(n>b/t\). For \(0<t<b\) and \(\sigma^2>0\),
\(B_1\leq B_2=\exp\{-(nt/b)h_1(\lambda)\}\) with \(\lambda=bt/\sigma^2\)
fixed, and \(h_1(\lambda)\geq h_2(\lambda)>0\) by
\cref{thm:p15-variance-tail-chain}; so both decrease geometrically in
\(n\). By contrast \(n\cdot\sigma^2/(\sigma^2+nt^2)\to\sigma^2/t^2\).
At \(n=1\), \(\sigma^2=1\), \(t=1/2\), the witness puts mass
\(1/(1+1/4)=4/5\) at \(1/2\) and mass \(1/5\) at \(-2\).
\end{proof}
\end{annexitem}

\begin{annexitem}{annex:p15-w02c-p15-017-nonpositive-deviation}{thm:p15-w02c-p15-017-nonpositive-deviation}{Constant and nearly constant witnesses}
\begin{proof}
\emph{Part 1.} Constant random variables are independent, since every
event they generate has probability zero or one. With \(X_i=p\in[0,1]\),
\(0\leq X_i\leq1\) and \(n^{-1}\sum_iEX_i=p\). Then \(\overline X-p=0\),
and \(0\geq t\), so \(\{\overline X-p\geq t\}\) is the whole space and has
probability one. Under the ranges of \cref{thm:p15-conc-unequal-ranges},
take \(X_i=c_i\) with \(a_i\leq c_i\leq b_i\); then
\(S_n=\sum_i(c_i-c_i)=0\geq t\), and again the event is certain.

\emph{Part 2.} Put \(s=\sigma\sqrt{nq/(1-q)}\) and
\(c=\sigma\sqrt{n(1-q)/q}\); both are positive since \(0<q<1\), the upper
limit \(b^2/(b^2+n\sigma^2)\) being smaller than one. Then
\[
 EX_1=(1-q)s-qc=\sigma\sqrt{nq(1-q)}-\sigma\sqrt{nq(1-q)}=0,
\]
\[
 EX_1^2=(1-q)s^2+qc^2=n\sigma^2q+n\sigma^2(1-q)=n\sigma^2 .
\]
The variables \(X_2,\ldots,X_n\) are zero, so the family is independent,
every mean is zero, and \(n^{-1}\sum_iEX_i^2=\sigma^2\). The largest value
taken is \(s\), and
\[
 s\leq b\iff\frac{n\sigma^2q}{1-q}\leq b^2\iff q\leq\frac{b^2}{b^2+n\sigma^2},
\]
so \(X_i\leq b\) for every \(i\), using \(b>0\) for the zero variables.
On \(\{X_1=s\}\), \(\overline X=s/n>0\geq t\). Hence
\(\Pr\{\overline X\geq t\}\geq\Pr\{X_1=s\}=1-q\). Letting
\(q\downarrow0\) shows that the supremum over the setting is at least
one; it is at most one because it is a supremum of probabilities.

\emph{Consequence.} A number \(\beta\) that bounds the probability for
every member of either class bounds it for the witnesses, so
\(\beta\geq1\) in part~1 and \(\beta\geq1-q\) for every admissible \(q\)
in part~2, whence \(\beta\geq1\).

\emph{The remark after the theorem.} At \(q=b^2/(b^2+n\sigma^2)\),
\(1-q=n\sigma^2/(b^2+n\sigma^2)\), so \(s=\sigma\sqrt{n\cdot b^2/(n\sigma^2)}=b\)
and \(c=\sigma\sqrt{n\cdot n\sigma^2/b^2}=n\sigma^2/b\). These are the
values \(b\) and \(-v/b\), with probabilities \(v/(b^2+v)\) and
\(b^2/(b^2+v)\), of \cref{thm:p15-upper-moment-extremal} for
\(v=n\sigma^2\).
\end{proof}
\end{annexitem}

\begin{annexitem}{annex:p15-w02c-p15-017-markov-comparison}{thm:p15-w02c-p15-017-markov-comparison}{Two elementary logarithm bounds}
\begin{proof}
Write \(q=p+t\), so \(p<q<1\). The ordering \(A_1\leq A_2\leq A_3\) is
\cref{thm:p15-conc-common-range}. Markov's inequality applies because
\(\overline X\geq0\) and \(E\overline X=p\):
\(\Pr\{\overline X\geq q\}\leq E\overline X/q=p/(p+t)\).

\emph{Part 1.} With \(n=1\), \(A_1=e^{-D(q\Vert p)}\), and
\[
 \ln\frac qp-D(q\Vert p)
 =(1-q)\ln\frac qp-(1-q)\ln\frac{1-q}{1-p}
 =(1-q)\ln\frac{q(1-p)}{p(1-q)} .
\]
Since \(q>p\), \(q(1-p)-p(1-q)=q-p>0\), so the logarithm is positive, and
\(1-q>0\). Hence \(D(q\Vert p)<\ln(q/p)\), that is,
\(A_1=e^{-D(q\Vert p)}>p/q=p/(p+t)\).

\emph{Part 2.} For \(x>0\), \(x\neq1\), \(\ln x<x-1\). Apply it to both
logarithms of \(D(q\Vert p)\); the ratios \(q/p\) and \((1-q)/(1-p)\)
differ from one:
\[
 D(q\Vert p)<q\Bigl(\frac qp-1\Bigr)+(1-q)\Bigl(\frac{1-q}{1-p}-1\Bigr)
 =(q-p)\Bigl(\frac qp-\frac{1-q}{1-p}\Bigr)
 =\frac{(q-p)^2}{p(1-p)} .
\]
The last step uses \(q(1-p)-p(1-q)=q-p\). So
\(A_1>\exp\{-nt^2/(p(1-p))\}\). Next, for \(x>0\),
\(\ln(1+x)>x/(1+x)\): the difference vanishes at \(x=0\) and has
derivative \(1/(1+x)-1/(1+x)^2=x/(1+x)^2>0\). With \(x=t/p\) this gives
\(\ln(1+t/p)>t/(p+t)\). If \(nt(p+t)\leq p(1-p)\), then
\[
 \frac{nt^2}{p(1-p)}\leq\frac t{p+t}<\ln\Bigl(1+\frac tp\Bigr),
\]
so \(A_1>\exp\{-\ln(1+t/p)\}=p/(p+t)\).

\emph{Part 3.} For \(x>0\), \(\ln(1+x)<x\): the difference
\(x-\ln(1+x)\) vanishes at zero and has derivative \(x/(1+x)>0\). If
\(2npt>1\), then \(2nt^2>t/p>\ln(1+t/p)\), so
\(A_3=e^{-2nt^2}<e^{-\ln(1+t/p)}=p/(p+t)\).

\emph{Part 4.} At \(t=1-p\), Markov's value is \(p/1=p\), and the
extension of \(A_1\) is \(p^n\) by \cref{thm:p15-conc-common-range}. For
\(n=1\) they are equal; for \(n\geq2\), \(p^n=p\cdot p^{n-1}<p\) since
\(0<p<1\).

\emph{The summary after the theorem.} The first clause is part~2. For
fixed \(p,t\), \(n>1/(2pt)\) is \(2npt>1\), which is part~3. The last
clause is part~1.

\emph{The example.} With \(p=1/2\), \(t=1/4\): \(p/(p+t)=2/3\);
\(g(1/2)=2\), so \(A_2=A_3=e^{-2n/16}=e^{-n/8}\);
\(D(3/4\Vert1/2)=\tfrac34\ln\tfrac32+\tfrac14\ln\tfrac12\), so
\(e^{-D}=(2/3)^{3/4}2^{1/4}=2/3^{3/4}\). The condition of part~2,
\(n\cdot\tfrac14\cdot\tfrac34\leq\tfrac14\), holds only for \(n=1\); that
of part~3, \(n/4>1\), holds for \(n\geq5\). Both \(A_1\) and \(A_3\)
decrease in \(n\). Now \(A_1<2/3\) is \(3\cdot2^{n-1}<3^{3n/4}\). At
\(n=3\) the left side is \(12\) and the right is
\(9\cdot3^{1/4}<9\cdot\tfrac43=12\), because \(3<(4/3)^4=256/81\); so
\(A_1>2/3\) for \(n\leq3\). At \(n=4\) the comparison is \(24<27\), so
\(A_1=16/27<2/3\). Also \(A_3<2/3\) is \(n>8\ln\tfrac32\). Here \(\ln\tfrac32<\tfrac12\) by
\(\ln(1+x)<x\), and \(\ln\tfrac32\geq\tfrac{2\cdot(1/2)}{2+1/2}=\tfrac25>\tfrac38\)
by \(\ln(1+x)\geq2x/(2+x)\), proved for \cref{thm:p15-variance-tail-chain};
so \(3<8\ln\tfrac32<4\), and \(A_3<2/3\) holds exactly for \(n\geq4\).
\end{proof}
\end{annexitem}

\begin{annexitem}{annex:p15-w02c-p15-017-section-three-reading}{thm:p15-w02c-p15-017-section-three-reading}{Bracketing a ratio}
\begin{proof}
\Cref{thm:p15-conc-common-comparison} gives \(L<A_1/A_2<U\) for
\(p>1/2\) and for \(p=1/2\). Each exponent in \(U\) is positive: for
\(p>1/2\) because \(2p-1>0\), and for \(p=1/2\) because \(t>0\). So
\(U<1\), and \(L>0\) as an exponential. Multiplying \(A_1/A_2<U\) by
\(A_2>0\) gives \(A_1<UA_2\). Also
\((A_2-A_1)/A_2=1-A_1/A_2\), which lies strictly between \(1-U>0\) and
\(1-L\).

For the factor between the exponents when \(p>1/2\), divide the exponent
of \(L\) by that of \(U\):
\[
 \frac{nt^3}{3(1-p)(1-p-t)}\cdot\frac{6p^2(1-p)^2}{(2p-1)nt^3}
 =\frac{2p^2(1-p)}{(2p-1)(1-p-t)} .
\]
Since \(0<1-p-t<1-p\), the denominator is less than \((2p-1)(1-p)\), and
\(2p^2-(2p-1)=p^2+(1-p)^2>0\); so the factor exceeds
\(2p^2(1-p)/\{(2p-1)(1-p)\}>1\). As \(p\downarrow1/2\) with \(t\) fixed
the factor tends to infinity, because \(2p-1\to0\) while the numerator
tends to \(1/4\) and \(1-p-t\) to \(1/2-t>0\).

For \(p<1/2\), the cited theorem gives
\(e^{-n(1-2p-t)^2/p}\leq A_1/A_2\leq1\), so
\(0\leq(A_2-A_1)/A_2\leq1-e^{-n(1-2p-t)^2/p}\).

\emph{The example.} With \(p=3/4\), \(t=1/8\): \(2p-1=1/2\),
\(6p^2(1-p)^2=6\cdot\tfrac9{16}\cdot\tfrac1{16}=\tfrac{27}{128}\) and
\(t^3=1/512\), so the exponent of \(U\) is
\(\tfrac12\cdot\tfrac n{512}\cdot\tfrac{128}{27}=\tfrac n{216}\). Also
\(3(1-p)(1-p-t)=3\cdot\tfrac14\cdot\tfrac18=\tfrac3{32}\), so the
exponent of \(L\) is \(\tfrac n{512}\cdot\tfrac{32}3=\tfrac n{48}\). At
\(n=4\), \(1-L=1-e^{-1/12}<1/12\) because \(e^{-x}>1-x\) for \(x>0\). At
\(n=432\), \(U=e^{-2}\), and \(e^{-2}<0.14\) because \(e^2>7.14\), which
follows from \(e^2>1+2+2+\tfrac43+\tfrac23+\tfrac4{15}=7.2\overline6\)
(the first six terms of the exponential series).
\end{proof}
\end{annexitem}

\begin{annexitem}{annex:p15.concentration-convex-mixture}{thm:p15-concentration-convex-mixture}{Center first, then apply finite Jensen}
\begin{proof}
Discard any zero-weight terms.
Linearity gives \(T-ET=\sum_jp_j(T_j-ET_j)\).
Finite Jensen for the convex function \(x\mapsto e^{hx}\) gives
the pointwise inequality
\[
 e^{h(T-ET)}\leq\sum_jp_je^{h(T_j-ET_j)}.
\]
All terms are nonnegative, so expectations preserve the inequality
even if some are infinite. The exponential indicator bound gives the
tail conclusion. No factorization across different \(T_j\) occurred.

For (5.2), fix real \(t\) and \(h>0\). Finite Jensen for
\(x\mapsto e^{hx}\) gives \(e^{hT}\leq\sum_jp_je^{hT_j}\) pointwise, and
\({\bf1}_{\{T\geq t\}}\leq e^{h(T-t)}\). Multiplying the first inequality
by \(e^{-ht}>0\) and taking expectations of nonnegative variables gives
\(\Pr\{T\geq t\}\leq\sum_jp_jEe^{h(T_j-t)}\). Replacing each \(T_j\) by
\(T_j-ET_j\), which replaces \(T\) by
\(\sum_jp_j(T_j-ET_j)=T-ET\), yields the centered display. If
\(ET_j=ET\) for all \(j\), then
\(Ee^{h(T_j-ET-t)}=Ee^{h(T_j-ET_j-t)}\) termwise, so (5.2) at the
threshold \(ET+t\) is the centered bound.

If a number \(\beta\) satisfies \(Ee^{h(T_j-t)}\leq\beta\) for every
\(j\), then (5.2) gives \(\Pr\{T\geq t\}\leq\beta\), while the
exponential indicator bound for the single variable \(T_j\) gives
\(\Pr\{T_j\geq t\}\leq Ee^{h(T_j-t)}\leq\beta\). Identically distributed
\(T_j\) have equal moments, so this case is included. Finally, the
weights are nonnegative with sum one, so
\(\sum_jp_jEe^{h(T_j-t)}\geq\min_jEe^{h(T_j-t)}\).
\end{proof}
\end{annexitem}

\begin{annexitem}{annex:p15.one-sample-u-concentration}{thm:p15-one-sample-u-concentration}{Counting ordered blocks and retaining their centering}
\begin{proof}
Fix an ordered distinct \(r\)-tuple. In each of the \(k\) block positions
it occurs in exactly \((n-r)!\) permutations.
Its coefficient in \(n!^{-1}\sum_\pi V_\pi\) is therefore
\(k(n-r)!/(kn!)=1/(n)_r\), proving the identity pointwise.
For a fixed permutation the \(k\) blocks are functions of disjoint
subfamilies of jointly independent inputs, hence mutually independent.
By \cref{thm:p15-conc-exponential}, their centered average satisfies
\[
 E e^{h(V_\pi-EV_\pi)}
 \leq\prod_{j=1}^k
       \exp\{h^2(b-a)^2/(8k^2)\}
 =e^{h^2(b-a)^2/(8k)}.
\]
Apply the mixture lemma to \(U-EU\), retaining \(EV_\pi\) separately
when the inputs have unequal laws. Minimize
\(-ht+h^2(b-a)^2/(8k)\) at \(h=4kt/(b-a)^2\); the value is
\(-2kt^2/(b-a)^2\).

With identically distributed inputs, every block has the same mean
\(\mu\) and the same kernel law. For a kernel in \([0,1]\), the chord
bound gives the common moment envelope
\[
 E e^{h(V_\pi-\mu)}
 \leq\{e^{-h\mu/k}(1-\mu+\mu e^{h/k})\}^{k}.
\]
The mixture lemma transfers this entire envelope. Its minimization
and both subsequent comparisons are precisely
\cref{thm:p15-conc-common-range} with count \(k\).
At an upper endpoint, increase thresholds from below, so the events
decrease to the endpoint event; beyond
the endpoint the statistic cannot exceed its kernel's upper bound.
For \(\mu=0,1\), boundedness forces each of the finitely many kernel
values in the statistic to be its mean almost surely.

For a centered upper-bounded kernel with finite variance, the common
envelopes instead are \(M_{b_0,\sigma_g^2}(h/k)^k\) and its tangent
envelope. Transfer them first and optimize afterward, exactly as in
\cref{thm:p15-variance-tail-chain}. The same zero-variance and
endpoint arguments apply. No independence between overlapping
kernel terms in \(U\) is assumed.
\end{proof}
\end{annexitem}

\begin{annexitem}{annex:p15.two-sample-u-concentration}{thm:p15-two-sample-u-concentration}{A product permutation count}
\begin{proof}
Fix an ordered \(r\)-tuple from the first sample and ordered \(s\)-tuple
from the second. In any of the \(k\) paired block positions this choice
occurs in \((m-r)!(n-s)!\) permutation pairs. Across all positions the
count is \(k(m-r)!(n-s)!\). After division by \(k m!n!\), its coefficient is
\(1/\{(m)_r(n)_s\}\), proving the identity.
Each paired block uses inputs disjoint from all other paired blocks.
Mutual independence of all inputs therefore gives independence of
these \(k\) kernel values.
The exponential bound for their centered average is
\(e^{h^2(b-a)^2/(8k)}\); mixture and optimization give the stated tail.
When each sample has a common law within itself, every paired block
has the same mixed kernel law, even if the two laws differ.
Its chord or upper-variance moment envelope is then the same in
every permutation pair. The preceding proof transfers each envelope,
its optimization, and its endpoint/zero-variance qualifications verbatim
with this specified \(k\).
\end{proof}
\end{annexitem}

\begin{annexitem}{annex:p15.repeated-index-concentration}{thm:p15-repeated-index-concentration}{Partitioning every repeated-index pattern}
\begin{proof}
Partition \(\{1,\ldots,r\}\) according to equality of the chosen indices
in an index tuple. A fixed partition \(\mathcal P\) with \(j\) blocks
has \((n)_j\) index tuples, one for each injective assignment of its
labeled blocks to population indices. Consequently
\[
 n^r=\sum_{\mathcal P}(n)_{|\mathcal P|},\qquad
 W=\sum_{\mathcal P}\frac{(n)_j}{n^r}U_j(g_{\mathcal P}),
\]
where
\(g_{\mathcal P}(x_1,\ldots,x_j)
 =g(x_{\beta_{\mathcal P}(1)},\ldots,x_{\beta_{\mathcal P}(r)})\).
For a kernel of the first \(j\) of \(r\) arguments, its ordered
\(U_r\) average equals its ordered \(U_j\) average: every distinct
\(j\)-tuple has \((n-j)_{r-j}\) completions, and
\((n)_r=(n)_j(n-j)_{r-j}\).
Apply this identity term by term to the displayed sum to obtain
\(W=U_r(g^*)\). The coefficients are nonnegative and sum to one,
so they preserve the bounds \(a,b\). The one-sample theorem applies.

For \(r=2\) there are the distinct and all-equal patterns, of weights
\((n)_2/n^2\) and \(n/n^2\). For \(r=3\) there are the all-distinct
pattern, three patterns with exactly two blocks, and the all-equal
pattern. Their weights are respectively
\((n)_3/n^3,(n)_2/n^3,n/n^3\).
Choose the block labels in the three middle patterns to produce
\(g(x_1,x_1,x_2),g(x_1,x_2,x_1),g(x_2,x_1,x_1)\).
Their displayed coefficients are exactly those in the body.
The arbitrary labeling is legitimate because each injective assignment
is averaged; no symmetry of \(g\) is needed.

For the placement statement, let \(f\) have \(j\leq r\) arguments, let
\(\iota\) be injective, and put
\(f_\iota(x_1,\ldots,x_r)=f(x_{\iota(1)},\ldots,x_{\iota(j)})\).
Send a distinct-index \(r\)-tuple \((i_1,\ldots,i_r)\) to the distinct
\(j\)-tuple \((i_{\iota(1)},\ldots,i_{\iota(j)})\). Every distinct
\(j\)-tuple is the image of exactly \((n-j)_{r-j}\) distinct \(r\)-tuples:
the \(r-j\) positions outside the image of \(\iota\) receive distinct
indices from the \(n-j\) unused ones, in order. Hence
\[
 \sum_{i_1,\ldots,i_r\ {\rm distinct}}f_\iota(X_{i_1},\ldots,X_{i_r})
 =(n-j)_{r-j}\sum_{i_1,\ldots,i_j\ {\rm distinct}}
 f(X_{i_1},\ldots,X_{i_j}),
\]
and dividing by \((n)_r=(n)_j(n-j)_{r-j}\) gives
\(U_r(f_\iota)=U_j(f)\). The case of the first \(j\) positions is the
identity used above. Since \(U_r\) is linear in its kernel, a convex
combination of placements of \(f\) also has ordered average \(U_j(f)\);
replacing a partition term of \(g^*\) by such a combination leaves
\(W=U_r(g^*)\) intact, and the new kernel is still a weighted mean of
values of \(g\), so it still lies in \([a,b]\).

For \(r=2\) the all-equal pattern is \(f(x)=g(x,x)\) with \(j=1\).
Its placements at position one and at position two are \(g(x_1,x_1)\)
and \(g(x_2,x_2)\), each with ordered average \(n^{-1}\sum_ig(X_i,X_i)\).
Splitting its weight \(n/n^2=1/n\) equally between them gives
\(\tilde g^*\), whose coefficients are nonnegative and sum to one.
If \(g(x,y)=g(y,x)\), interchanging \(x_1,x_2\) maps each of the three
terms of \(\tilde g^*\) to a term with the same coefficient. For the
same symmetric \(g\), the kernel (5.16) satisfies
\(g^*(x_1,x_2)-g^*(x_2,x_1)=\{g(x_1,x_1)-g(x_2,x_2)\}/n\), so it is
symmetric exactly when \(x\mapsto g(x,x)\) is constant.
\end{proof}
\end{annexitem}

\begin{annexitem}{annex:p15.cramer-von-mises-kernel}{thm:p15-cramer-von-mises-kernel}{Expanding the empirical-distribution integral}
\begin{proof}
Continuity of \(G\) means its probability measure has no atoms.
For fixed \(x\), with \(u=G(x)\), integration with respect to \(dG\)
gives
\[
 \int{\bf1}_{\{x\leq z\}}\,dG(z)=1-u,\quad
 \int{\bf1}_{\{x\leq z\}}G(z)\,dG(z)=\frac{1-u^2}{2},\quad
 \int G(z)^2\,dG(z)=\frac13.
\]
One elementary justification of the substitution used here is to
partition a finite interval and compare the Stieltjes sums
\(\sum G(z_j)^p\{G(z_j)-G(z_{j-1})\}\) with the ordinary Riemann
sums for \(\int s^pds\) over the corresponding \(G\)-increments.
Continuity makes the maximum increment tend to zero; boundedness
then gives the integral identity on that interval. Let its endpoints
tend to infinity, using \(G(-\infty)=0,G(+\infty)=1\), and use the
same argument from a fixed \(x\). This proves the displayed identities
also for continuous distributions with flat portions.

Expand the square defining \(\omega^2\). The \(F_n^2\) term is
\(n^{-2}\sum_{i,j}\{1-\max(G(X_i),G(X_j))\}\).
The \(-2F_nG\) term is
\(-n^{-1}\sum_i\{1-G(X_i)^2\}\), and the remaining term is \(1/3\).
Distribute the single-index sum equally across ordered pairs to get
the stated kernel.

For \(0\leq u\leq v\leq1\), the kernel is
\(1/3+u^2/2+v^2/2-v\).
It is minimized in \(u\) at zero, and the remaining decreasing function
of \(v\) is minimized at one, giving \(-1/6\).
For its maximum with \(v\) fixed take \(u=v\), giving
\(1/3+v^2-v\leq1/3\). Interchange \(u,v\) for the other triangle.
The kernel width is \(1/2\), so the repeated-index theorem gives
the exponent \(-2\lfloor n/2\rfloor t^2/(1/2)^2\).
The identity was pointwise and the concentration premise only required
independent \(X_i\); no null hypothesis \(X_i\sim G\) was used.
At \(u=1/4,v=3/4\), the two diagonal kernel values are \(7/48\)
and both off-diagonal values are \(-5/48\).
Their four-term average is \(1/48\).

For the range of \(g^*\), put \(\varepsilon=1/n\in(0,1/2]\),
\(u=G(x_1)\), \(v=G(x_2)\). Then
\[
 g^*=\frac13+(1-\varepsilon)\Bigl(\frac{u^2}2+\frac{v^2}2-\max\{u,v\}\Bigr)
 +\varepsilon(u^2-u).
\]
Upper bound: \(g\leq1/3\) was shown above, and
\(g(x_1,x_1)=1/3+u^2-u\leq1/3\) since \(u^2\leq u\); the convex
combination is at most \(1/3\).
Lower bound, case \(u\leq v\): on \(u\leq v\leq1\) the derivative in
\(v\) is \((1-\varepsilon)(v-1)\leq0\), so \(g^*\) is at least its value
at \(v=1\),
\[
 \frac13-\frac{1-\varepsilon}2+\frac{1+\varepsilon}2u^2-\varepsilon u
 \geq\frac13-\frac{1-\varepsilon}2-\frac{\varepsilon^2}{2(1+\varepsilon)}
 =\frac{2\varepsilon-1}{6(1+\varepsilon)},
\]
the quadratic in \(u\) being minimized at
\(u=\varepsilon/(1+\varepsilon)\in[0,1]\).
Case \(v\leq u\): the only \(v\)-dependent term is
\((1-\varepsilon)v^2/2\geq0\), so \(g^*\) is at least its value at
\(v=0\),
\[
 \frac13+\frac{1+\varepsilon}2u^2-u\geq\frac13-\frac1{2(1+\varepsilon)}
 =\frac{2\varepsilon-1}{6(1+\varepsilon)},
\]
minimized at \(u=1/(1+\varepsilon)\in(0,1)\).
With \(\varepsilon=1/n\) the common lower bound is \((2-n)/\{6(n+1)\}\).

Both bounds are approached. A continuous distribution function tends to
\(0\) and \(1\) at \(-\infty\) and \(+\infty\), so by the intermediate
value theorem it takes every value in \((0,1)\). Choose \(x_1\) with
\(G(x_1)=n/(n+1)\) and let \(x_2\to-\infty\); then \(v\to0\) and, by
continuity of the displayed polynomial expression in \((u,v)\),
\(g^*\) tends to the lower bound. Letting \(x_1=x_2\to-\infty\) sends
\(g^*\) to \(1/3\). For \(g\) itself, \(u\to0,v\to1\) gives \(-1/6\)
and \(u=v\to0\) gives \(1/3\). Hence the width of \(g^*\) is
\[
 \frac13-\frac{2-n}{6(n+1)}=\frac{3n}{6(n+1)}=\frac n{2(n+1)}<\frac12.
\]
By \cref{thm:p15-repeated-index-concentration}, \(\omega^2=U_2(g^*)\)
for \(n\geq2\), and the one-sample theorem with this width gives the
exponent \(-2\lfloor n/2\rfloor t^2\{2(n+1)/n\}^2
=-8\lfloor n/2\rfloor(1+1/n)^2t^2\).
Relative to the exponent from \(g\), it is larger by the factor
\((1+1/n)^2=1+O(1/n)\). At \(n=2\) the lower bound is \(0\).
\end{proof}
\end{annexitem}

\begin{annexitem}{annex:p15.finite-range-concentration}{thm:p15-finite-range-concentration}{Joint independence within residues and the unequal-size step}
\begin{proof}
In a fixed residue class, append indices in increasing order.
The next index is at least \(r\) beyond the preceding one.
The dependence premise makes its tail vector independent of the
entire prefix containing all earlier indices in the residue.
Induction on the number of selected indices therefore factors every
finite intersection of their coordinate events, proving mutual,
not merely pairwise, independence within the class.
The residue sets partition \(\{1,\ldots,n\}\), each \(n_i\geq k\geq1\),
and \(\sum_i n_i=n\). Summing the centered residue averages with
weights \(n_i/n\) proves the displayed decomposition.
The convex-mixture lemma gives its exact exponential inequality.

For bounded \(Y_j\), independence within a residue and the centered
exponential lemma give
\[
 E e^{h(S_i-ES_i)/n_i}
 \leq e^{h^2(b-a)^2/(8n_i)}
 \leq e^{h^2(b-a)^2/(8k)}.
\]
Transfer this common bound and optimize in \(h\) as before.
For common distributions, let \(L(h)\) be the log moment envelope
of one centered Bernoulli endpoint law for the common-range case,
or \(L(h)=\log M_{b_0,\sigma^2}(h)\) in the variance case.
Each is the log moment function of a bounded two-point centered law;
\(L(0)=0\) and \(L\) is convex. Explicitly,
its second derivative is the variance under the exponentially tilted
two-point probabilities and is nonnegative.
For integers \(q\geq k\), convexity at
\(h/q=(k/q)(h/k)+(1-k/q)0\) gives
\[
 qL(h/q)\leq kL(h/k).
\]
Hence the entire residue envelope \(e^{n_iL(h/n_i)}\) is at most
\(e^{kL(h/k)}\). This supplies the missing common bound when the
sizes are unequal. Applying the mixture lemma and then the earlier
optimizations proves all claimed common-range and variance bounds.
Endpoint extensions follow by increasing thresholds from below and
continuity for the resulting decreasing events. Zero
variance makes every individual variable constant almost surely.
When \(r>n\), replacing \(r\) by \(n\) leaves a vacuous separation
premise, gives singleton residues and \(k=1\), so no division by an
empty class is introduced.
\end{proof}
\end{annexitem}

\begin{annexitem}{annex:p15.without-replacement-convex}{thm:p15-without-replacement-convex}{An explicit symmetrized version of the source's comparison}
\begin{proof}
All averages here are finite. For a partition \(\mathcal P\) of
\(\{1,\ldots,n\}\) into \(j\) blocks, label the blocks and let their
sizes be \(r_1,\ldots,r_j\), so \(\sum_\ell r_\ell=n\).
Put \(p_{\mathcal P}=(N)_j/N^n\). Partitioning all \(N^n\) possible
with-replacement index tuples by their equality patterns gives
\(\sum_{\mathcal P}p_{\mathcal P}=1\).
Define, for a deterministic \(n\)-tuple \(x\),
\[
 \overline g(x_1,\ldots,x_n)=
 \sum_{\mathcal P}p_{\mathcal P}\frac1{(n)_j}
 \sum_{i_1,\ldots,i_j\ {\rm distinct}}
 f\!\left(\sum_{\ell=1}^j r_\ell x_{i_\ell}\right).
\]
This provides explicit nonnegative coefficients for the source's
symmetrized repeated-index kernel. For a without-replacement sample,
any fixed ordered \(j\) of its coordinates is uniform over distinct
population indices: each such choice has \((N-j)_{n-j}\) completions
among the \((N)_n\) sample tuples.
Consequently, taking expectation in the last display and grouping
with-replacement tuples by \(\mathcal P\) gives exactly
\[
 E\overline g(X_1,\ldots,X_n)=E f(Y_1+\cdots+Y_n).
\]

For each fixed partition, the average over injective \(i_1,\ldots,i_j\)
places a uniform coordinate in every one of its \(j\) positions.
The average of its arguments of \(f\) is therefore
\[
 \sum_{\ell=1}^j r_\ell\frac1n\sum_{i=1}^nx_i
 =\sum_{i=1}^nx_i.
\]
All partition weights sum to one, so the same holds for the total
weighted average. Finite Jensen gives the pointwise inequality
\(\overline g(x)\geq f(\sum_i x_i)\).
Taking expectation proves the convex comparison.
This calculation proves the source's required coefficient and
linear-mean identities without assuming an independent law for \(X\).
Finiteness and continuity of \(f\) make all its values at the finitely
many possible sums finite.

Set \(f(s)=e^{hs}\) and multiply by \(e^{-hn\mu}\).
The centered without-replacement exponential moment is bounded by
that of \(n\) independent uniform population draws. Every tail
bound obtained from that independent exponential envelope therefore
transfers, with exactly its independent-law parameters.
In particular \(\sigma_C^2\), not a reduced variance substituted by
analogy, is the parameter in the proved variance envelope.

For \(N>1\), set \(d_j=c_j-\mu\). For two distinct sampled coordinates,
\[
 E[(X_1-\mu)(X_2-\mu)]
 =\frac{\sum_{i\ne j}d_id_j}{N(N-1)}
 =-\frac{\sum_jd_j^2}{N(N-1)}
 =-\frac{\sigma_C^2}{N-1}.
\]
Each coordinate has variance \(\sigma_C^2\). Thus
\[
 \operatorname{Var}\Big(\sum_iX_i\Big)
 =n\sigma_C^2-n(n-1)\frac{\sigma_C^2}{N-1}
 =n\sigma_C^2\frac{N-n}{N-1}.
\]
Divide by \(n^2\). For independent draws the cross-covariances are
zero, giving the other formula. Each coordinate has mean \(\mu\).
For \(N=1\), all coordinates are deterministic.
At \(n=N\), the without-replacement sum is the fixed population sum,
establishing the stated zero-deviation boundary.

The function \(f(s)=(s-n\mu)^2\) is continuous and convex. Both sums
have mean \(n\mu\), since each coordinate has mean \(\mu\). Hence
\(Ef(\sum_iX_i)=\operatorname{Var}(\sum_iX_i)\) and
\(Ef(\sum_iY_i)=\operatorname{Var}(\sum_iY_i)\), and the convex comparison
gives \(\operatorname{Var}(\sum_iX_i)\leq\operatorname{Var}(\sum_iY_i)\).
Dividing by \(n^2\) gives
\(\operatorname{Var}(\overline X)\leq\operatorname{Var}(\overline Y)\),
in agreement with the explicit formulas, since \((N-n)/(N-1)\leq1\).

For \(n=2\) there are two partitions of \(\{1,2\}\). The partition into
singletons has \(j=2\), sizes \(r_1=r_2=1\), and weight
\((N)_2/N^2=(N-1)/N\); its inner average runs over the ordered pairs
\((1,2),(2,1)\) and equals \(f(x_1+x_2)\). The one-block partition has
\(j=1\), size \(r_1=2\), and weight \((N)_1/N^2=1/N\); its inner average
runs over \(i_1\in\{1,2\}\) and equals
\(\frac12\{f(2x_1)+f(2x_2)\}\). The general formula therefore gives
\[
 \overline g(x_1,x_2)=\frac{N-1}Nf(x_1+x_2)
 +\frac1{2N}f(2x_1)+\frac1{2N}f(2x_2).
\]
The weights sum to one and the weighted mean of the arguments is
\(\frac{N-1}N(x_1+x_2)+\frac1{2N}(2x_1+2x_2)=x_1+x_2\), so Jensen gives
\(\overline g(x_1,x_2)\geq f(x_1+x_2)\); the expectation identity is the
case \(n=2\) of the one proved above.
\end{proof}
\end{annexitem}

\begin{annexitem}{annex:p15.dependent-concentration-examples}{thm:p15-dependent-concentration-examples}{Exact values and the dependence countermodel}
\begin{proof}
For the variance example, let \(K\) count the variables equal to one.
It has the binomial law with count twenty and success probability \(1/5\).
Each summand has mean \(1/5-(4/5)(1/4)=0\) and second moment
\(1/5+(4/5)(1/16)=1/4\). The sum is \(5K/4-5\).
Thus its mean is at least \(1/4\) exactly when \(K\geq8\), giving
the displayed finite binomial sum. Here \(\lambda=1,\tau=5\).
In \(B_1\) the two exponents before multiplying by twenty are
\(-2/5,-3/5\), so
\(B_1=2^{-8}(3/4)^{-12}=2^{16}/3^{12}\).
The other three values follow from
\(h_1(1)=2\log2-1\), \(h_2(1)=3/8\), and
\(h_3(1)=\operatorname{arsinh}(1/2)/2\).
Evaluating these explicit finite expressions gives the stated decimals.

For the unbounded-below example, geometric sums give
\[
 \sum_{j\geq1}7\,8^{-j}=1,\quad
 EZ=7\sum_{j\geq1}4^{-j}=7/3,\quad
 EZ^2=7\sum_{j\geq1}2^{-j}=7.
\]
Hence \(EX=0\), \(EX^2=7-49/9=14/9\), and the smallest value
of \(Z\) is two, so \(X\leq1/3\).
For \(h>0\), the summands \(7\,8^{-j}e^{h2^j}\) do not even tend
to zero, since \(h2^j-j\log8\to\infty\).
Their sum is infinite; multiplication by \(e^{-7h/3}\) proves the
assertion about \(Ee^{-hX}\).

The sign kernel lies in \([-1,1]\). Substitution of width two and
order two into the one-sample result gives the stated bound.
For the three listed observations, the unordered pairs have signs
\(+1,+1,-1\); reversing a pair leaves its sign unchanged.
The ordered sum is therefore two and its six-term average is \(1/3\).
This definition permits ties through sign zero; continuity of the
coordinate distributions is an additional way of excluding marginal
ties, not needed for the bound.

For the strict-comparison kernel the range is \([0,1]\) and each
paired block contains one first-sample and one second-sample input.
Thus the two-sample theorem gives \(k=\min(m,n)\) and the displayed
bound. In the listed arrays, \(X=1\) exceeds just \(Y=0\);
\(X=3\) exceeds \(Y=0,2\). The count is three of six.
This is the normalized pair count described in the source, with the
strict inequality retained.

For the moving-product example, the vector \((Y_1,\ldots,Y_i)\)
is measurable in \(Z_1,\ldots,Z_{i+r-1}\), while
\((Y_j,\ldots,Y_n)\) is measurable in
\(Z_j,\ldots,Z_{n+r-1}\).
When \(j-i\geq r\), these are disjoint independent input subfamilies,
so the required vectors are independent.
For Bernoulli inputs, a product is one exactly when all its \(r\)
factors are one, with probability \(p^r\). This gives its common
mean, range and distribution. At \(n=5,r=2\), the residues are
\(\{1,3,5\}\) and \(\{2,4\}\); the minimum size is two.
Substitution into the range theorem gives \(e^{-4t^2}\).

For the source's product example, \((Y_1,\ldots,Y_i)\) is measurable in
\(X_1,\ldots,X_{i+r-1}\) and \((Y_j,\ldots,Y_n)\) in
\(X_j,\ldots,X_{n+r-1}\). When \(j-i\geq r\), the first family ends at
index \(i+r-1<j\), so the two input families are disjoint and,
by independence of the \(X_i\), the two vectors are independent.
At \(r=2\), \(Y_i=X_iX_{i+1}\), the moving product with \(Z=X\).
For \(r\geq3\), the factor indices \(\{i,i+r-1\}\) and
\(\{i+1,i+r\}\) are disjoint, as \(i+r-1\neq i+1\) and
\(i\neq i+r\).
For the lower range, let \(2\leq r\leq n\). The indices \(1,r,2r-1\)
are distinct and at most \(n+r-1\), so with Bernoulli\((p)\) inputs
\(\Pr\{Y_1=1,Y_r=1\}=\Pr\{X_1=X_r=X_{2r-1}=1\}=p^3\), while
\(\Pr\{Y_1=1\}=\Pr\{Y_r=1\}=p^2\). For \(0<p<1\), \(p^3\neq p^4\).
Independence of \(Y_1\) and \(Y_r\) is what \((r-2)\)-dependence would
require for the vectors \((Y_1)\) and \((Y_r,\ldots,Y_n)\), whose index
gap is \(r-1\); it fails.
For \(r\geq2\), the two factors of each \(Y_i\) are distinct independent
Bernoulli\((p)\) variables, so \(Y_i\in\{0,1\}\) with
\(\Pr\{Y_i=1\}=p^2\); all \(Y_i\) share this law and \(ES=np^2\).
\Cref{thm:p15-finite-range-concentration} with \(a=0,b=1\) and
\(k=\lfloor n/r\rfloor\) gives the displayed bound.

In the population example, there are \(\binom42=6\) unordered
distinct-index samples. One selects both zeros, four select one of
each value, and one selects both ones. With replacement, two independent
fair binary values instead give binomial probabilities \(1/4,1/2,1/4\).
Both mean values are \(1/2\) and \(\sigma_C^2=1/4\);
the variance formulas give \((2/3)(1/8)=1/12\) and \(1/8\).
Their exponential moments are
\((1+4e^h+e^{2h})/6\) and \((1+2e^h+e^{2h})/4\);
subtracting yields \((e^h-1)^2/12\).
The upper endpoint of the sample mean selects two ones, with exact
without-replacement probability \(1/6\).
The common-range endpoint bound is \((1/2)^2=1/4\);
the variance endpoint gives the same number using
\(b=1/2,\sigma^2=1/4,n=2\).

Finally the possible triples \((A,B,AB)\), each of probability \(1/4\),
are
\[
 (1,1,1),\quad(1,-1,-1),\quad(-1,1,-1),\quad(-1,-1,1).
\]
Every coordinate is a fair sign. Every pair attains its four sign
combinations once, so each pair is independent.
All three signs are positive in exactly one of the four outcomes.
Their mean therefore equals one with probability \(1/4\).
The independent theorem at \(n=3,b=1,\sigma^2=1,t=b\) would give
\(\{1/(1+1)\}^3=1/8\), a false conclusion for this joint law.
The failed premise is mutual independence.
\end{proof}
\end{annexitem}
